\documentclass[minitoc, hidelinks, headlogo, framedcover]{tfeEPCC}
\usepackage[spanish, es-tabla]{babel} % Comentar si se va a realizar en inglés
%\usepackage{parskip} % Quitar comentario si se quieren distinguir los párrafos con saltos y quitar las sangrías

\usepackage{tabularx}
\usepackage{booktabs}

% Rellenar con los datos del estudiante y el trabajo
\estudiante{[Nombre del autor]}{Grado en Ingeniería Informática en Ingeniería del Software}{Mapeo de SpecKit contra flujos profesionales de ingeniería de requisitos basados en normas}
\title{Mapeo de SpecKit contra flujos profesionales de ingeniería de requisitos basados en normas}
\keywords{Ingeniería de requisitos\and IREB\and SpecKit\and Spec-Driven Development\and ISO 29148\and mapeo de procesos}

\begin{document}
\sloppy 
\pagebreak 

%%%%%%%%%%%%% PORTADA %%%%%%%%%%%%%
% Poner convocatoria y año de presentación
\portadaTFE{Convocatoria, Año}

\pagestyle{empty}


%%%%%%%%%% CONTRAPORTADA %%%%%%%%%%
% Campos, separar con comas si se ponen varios: {
% studentGender
% tutor (obligatorio)
% tutorGender
% cotutor
% cotutorGender }
% Especificar un co-tutor lo añade a la contraportada
% Poner la letra f en cualquier campo de género añade una 'a' al título correspondiente (por ejemplo, studentGender = f hace que se muestre Autora)
\contraportadaTFE{tutor = Tutor}

%%%% CAPÍTULOS INTRODUCTORIOS %%%%
\pagenumbering{gobble}
%\chapter*{Agradecimientos} % Apartado opcional, descomentar si se quiere añadir
% Dedicatorias del trabajo, pueden ser en español o inglés dependiendo de lo que se considere más apropiado

\chapter*{Resumen}
Este trabajo presenta un mapeo sistem\'atico del flujo de trabajo de \textsc{SpecKit}, un \emph{toolkit} de c\'odigo abierto para \emph{Spec-Driven Development} (SDD) desarrollado por GitHub, contra un flujo profesional de ingenier\'ia de requisitos basado en normas (IREB CPRE Foundation Level, ISO/IEC/IEEE 29148 e INCOSE Guide for Writing Requirements). Se documenta el flujo real de \textsc{SpecKit} a partir de sus fuentes oficiales y se contrasta con las actividades fundamentales de RE. Para fundamentar el mapeo con evidencia emp\'irica, se construye un corpus de 52 requisitos extra\'idos de \emph{changelogs} de seis proyectos \emph{Open Source}, formalizados con atributos IREB. El resultado principal es un mapa de alineaci\'on que identifica la cobertura de cada actividad y criterio normativo por parte de \textsc{SpecKit}.

\noindent\textbf{Palabras clave}: \printkeywords.

\chapter*{Abstract}
This work presents a systematic mapping of \textsc{SpecKit}'s workflow ---an open-source toolkit for Spec-Driven Development (SDD) developed by GitHub--- against a professional norm-based requirements engineering workflow (IREB CPRE Foundation Level, ISO/IEC/IEEE 29148, and INCOSE Guide for Writing Requirements). The actual \textsc{SpecKit} flow is documented from official sources and contrasted with fundamental RE activities. To support the mapping with empirical evidence, a corpus of 52 requirements is built from \emph{changelogs} of six \emph{Open Source} projects, formalized with IREB attributes. The main result is an alignment map identifying the coverage of each normative activity and criterion by \textsc{SpecKit}.

\noindent\textbf{Keywords}: Requirements engineering, IREB, SpecKit, Spec-Driven Development, ISO 29148, process mapping.

% Índices con enlaces de las secciones, tablas y figuras
\newpage
\pagestyle{plain}
\pagenumbering{Roman}
\setcounter{page}{1}
\tableofcontents
%\pagebreak
\newpage
\listoftables
\newpage
\listoffigures
\newpage


%%%%%%% ESTRUCTURA DEL TFG %%%%%%%%%%%%%%
% A.PORTADA (según la estructura indicada a continuación) 
% B.CONTRAPORTADA (según la estructura indicada a continuación) 
% C.ÍNDICE GENERAL DE CONTENIDOS  
% D.ÍNDICE DE TABLAS  
% E.ÍNDICE DE FIGURAS 
% F.RESUMEN (Podrá incluirse también en inglés, si así lo indica el Tutor1) 
% G.CUERPO DEL TRABAJO (según la estr
% uctura indicada a continuación) 
% H.REFERENCIAS BIBLIOGRÁFICAS 
% (Según norma ISO690)
% I.ANEXOS, si los hubiera
%%%%%%%%%%%%%%%%%%%%%%%%%%%%%%%%%%%%%%%
%%%%%% Cuerpo del trabajo
% 1.INTRODUCCIÓN 
% 2.OBJETIVOS 
% 3.ANTECEDENTES / ESTADO DEL ARTE 
% 4.MÉTODOLOGÍA 
% 5.IMPLEMENTACIÓN Y DESARROLLO (Cuando proceda) 
% 6.RESULTADOS Y DISCUSIÓN 
% 7.CONCLUSIONES 
%%%%%%%%%%%%%%%%%%%%%%%%%%%%%%%%%%%%%%%%%%%
%%%%%%%%%%%%%%%%%%%%%%%%%%%%%%%%%%%%%%%%%%%
% Tamaño: Normalizado UNE A-4, salvo planos. 
% Tipo y tamaño de letra del texto: Times New Roman 12 pt, Arial 12 pt o similar. 
% Interlineado del texto: 1,5 líneas. 
% Márgenes del texto: Superior, Inferior y Derecha, 2,5 cm; Izquierda, 4 cm. 
% Numeración de páginas en margen inferior derecha y tamaño 8 pt. 
% Las  figuras  serán  numeradas  y  tituladas  debajo  de  las  mismas  (indicando  su  fuente  si  no  son  de  elaboración  
% propia). 
% Las  tablas  serán  numeradas  y  tituladas  encima  de  las  mismas  (indicando  su  fuente  si  no  son  de  elaboración  
% propia). 
%%%%%%%%%%%%%%%%%%%%%%%%%%%%%%%%%%%%%%%%%%%
%%%%%%%%%%%%%%%%%%%%%%%%%%%%%%%%%%%%%%%%%%%
%%%%%%%%%%%%%%%%%%%%%%%%%%%%%%%%%%%%%%%%%%%
%%%%%%%%%%%%%%%%%%%%%%%%%%%%%%%%%%%%%%%%%%%


\pagenumbering{arabic}
\pagestyle{fancy}
\setcounter{page}{1}

%%%%%%%% INTRODUCCIÓN %%%%%%%%
\import{chapters}{01_introduccion.tex}

%%%%%%%% OBJETIVOS %%%%%%%%
\import{chapters}{02_objetivos.tex}

%%%%%%%% ESTADO DEL ARTE %%%%%%%%
\import{chapters}{03_estado_del_arte.tex}

%%%%%%%% METODOLOGÍA %%%%%%%%
\import{chapters}{04_metodologia.tex}

%%%%%%%% IMPLEMENTACIÓN Y DESARROLLO %%%%%%%%
\import{chapters}{05_implementacion_y_desarrollo.tex}

%%%%%%%% RESULTADOS %%%%%%%%
\import{chapters}{06_resultados.tex}

%%%%%%%% CONCLUSIONES Y TRABAJOS FUTUROS %%%%%%%%
\import{chapters}{07_conclusiones_y_trabajos_futuros.tex}


% Si no se van a incorporar anexos, comentar el comando '\startappendices' junto a todos los imports de la carpeta 'appendices'
%%%%%%%%%%%%%%%%%%%%%%%%%%%%%%%%%%
%%%%%%%%%%% ANEXOS %%%%%%%%%%%%%%%
%%%%%%%%%%%%%%%%%%%%%%%%%%%%%%%%%%
\startappendices

%%%%%%%% ANEXO 1 %%%%%%%%
\import{appendices}{01_ejemplo_de_anexo.tex}

\pagebreak
%%%%%%%%%%%%%%%%%%%%%%%%%%%%%%%%%%
%%%%%%%%%% AL FINAL %%%%%%%%%%%%%%
%%%%%%%%%%%%%%%%%%%%%%%%%%%%%%%%%%
\pagestyle{empty}
%%%% https://en.wikibooks.org/wiki/LaTeX/Bibliography_Management
%%%% https://www.bibtex.com
%%%%%%%% BIBLIOGRAFÍA %%%%%%%%
% Se puede cambiar el estilo de la bibliografía, por defecto es 'unsrturl'
% Por ejemplo: \makebibliography{alphaurl}
\makebibliography
\end{document}